% Generated draft: requires mathematical and source review.
\documentclass[11pt]{article}
\usepackage[margin=25mm]{geometry}
\usepackage{fontspec}
\setmainfont{Latin Modern Roman}
\usepackage{amsmath,amssymb,mathtools,unicode-math}
\setmathfont{Latin Modern Math}
\usepackage{enumitem,needspace,xurl,hyperref}
\setlength{\parindent}{0pt}
\setlength{\parskip}{5pt}
\setlength{\emergencystretch}{3em}
\newcommand{\N}{\mathbb{N}}
\newcommand{\Z}{\mathbb{Z}}
\newcommand{\Q}{\mathbb{Q}}
\newcommand{\R}{\mathbb{R}}
\newcommand{\C}{\mathbb{C}}
\newcommand{\F}{\mathbb{F}}
\newcommand{\A}{\mathbb{A}}
\newcommand{\PP}{\mathbb{P}}
\newcommand{\E}{\mathbb{E}}
\newcommand{\eps}{\varepsilon}
\newcommand{\dd}{\,\mathrm{d}}
\begin{document}
\section*{UnsolvedMath ID: 8800089}
\textbf{Source problem code:} catalogue entry
\subsection{Definitions and mathematical statement}
For an imaginary quadratic order \(\mathcal O\), its class number is \(h(\mathcal O)=\#\operatorname{Pic}(\mathcal O)\); the odd-class-number condition means \(h(\mathcal O)\equiv1\pmod2\). If maximal orders are intended, this is the ideal class number of the associated imaginary quadratic field. A Heegner point on a parametrized elliptic curve is obtained by mapping an appropriate complex multiplication point of a modular curve to the elliptic curve, with the required discriminant and level conditions.

Rosen and Silverman ask whether the odd-class-number condition in the \emph{stated lower bound for Heegner points} is necessary. The lower bound and its full hypotheses are not supplied. The intended mathematical question is whether that particular conclusion remains valid with the parity hypothesis removed while keeping the other hypotheses, or whether an admissible even-class-number example violates it. This is distinct from asking whether a particular proof method uses oddness.

The excerpt does not identify the relevant field or order, elliptic curve, level, discriminant range, or quantity bounded. In particular a height bound, a degree bound, and a bound concerning independence or rank of Heegner points are different statements; none is selected here. The supplied summary mentions an additional no-complex-multiplication hypothesis in related results, but does not provide the theorem to which it belongs. The necessity question therefore cannot be completed into an exact numerical inequality without the missing slide context.
\subsection{Short English statement}
Can the specified Heegner-point lower bound retain its conclusion when the odd-class-number hypothesis is dropped? The bound itself and its other hypotheses are absent, so a precise parity-removal theorem or counterexample target cannot yet be fixed.
\subsection{Sources}
\begin{itemize}
\item Local UnsolvedMath registry, permanent ID 8800089, source code catalogue entry; attributed to Michael Rosen and Joseph Silverman; no proposal year supplied; update 2026-08-06.
\item IPAM, \emph{Workshop on Number Theory and Cryptography --- Open Problems} (2006), attachment 6669, PDF page 22, printed slide 16: \url{https://www.ipam.ucla.edu/programs/workshops/workshop-i-number-theory-and-cryptography-open-problems/?tab=schedule}. Registry access date 2026-07-29; PDF and legacy-PPT visual audit. Provenance reports 19 attachments, 17 PDFs and two legacy PPT decks, with repetitions and solved definitions excluded. Present status and redistribution terms require review.
\item The supplied summary mentions Rosen--Silverman-type Heegner results and no-CM and odd-class-number hypotheses without an exact lower bound or full reference. Its unresolved-status claim was not independently checked.
\end{itemize}
\end{document}
